\documentclass[UTF8, 12pt, fontset=fandol, oneside]{ctexart}
\usepackage{researchnotes}

\title{第十章\quad 双线性型}
\author{}
\date{}

\begin{document}

\maketitle

\section*{§10.1\quad 对偶空间}

设 $V$ 是数域 $\mathbb K$ 上的线性空间，称 $V\to \mathbb K$ 的线性映射（$\mathbb K$ 作为一维空间）为 $V$ 上的线性函数. 令 $V^*$ 为由 $V$ 上的线性函数全体组成的集合. 我们可以在 $V^*$ 上定义加法与数乘，使 $V^*$ 成为 $\mathbb K$ 上的线性空间（参考 §4.2）. $V^*$ 称为 $V$ 的共轭空间，当 $V$ 是有限维空间时，常称 $V^*$ 是 $V$ 的对偶空间.

现设 $V$ 是 $\mathbb K$ 上的 $n$ 维线性空间，$\{e_1,e_2,\cdots,e_n\}$ 是 $V$ 的一组基，则 $V\to \mathbb K$ 的任一线性函数 $f$ 被它在 $V$ 基上的值唯一确定. 现定义 $f_i(i=1,2,\cdots,n)$ 如下：
\[
  f_i(e_j)=\delta_{ij},
\]
这里 $\delta_{ij}$ 称为 Kronecker 符号，即当 $i=j$ 时，$\delta_{ii}=1$；当 $i\ne j$ 时，$\delta_{ij}=0$. 我们得到了 $n$ 个 $V$ 上的线性函数 $\{f_1,f_2,\cdots,f_n\}$，现要证明这 $n$ 个线性函数组成了 $V^*$ 的一组基. 注意若
\[
  x=a_1e_1+a_2e_2+\cdots+a_ne_n,
\]
则
\[
  f_i(x)=a_i,\quad i=1,2,\cdots,n.
\]
设 $f$ 是 $V^*$ 中任一元素，且 $f(e_j)=b_j(j=1,2,\cdots,n)$，令
\[
  g=b_1f_1+b_2f_2+\cdots+b_nf_n,
\]
则
\[
\begin{aligned}
  g(x)&=(b_1f_1+b_2f_2+\cdots+b_nf_n)(x)\\
      &=b_1a_1+b_2a_2+\cdots+b_na_n\\
      &=a_1f(e_1)+a_2f(e_2)+\cdots+a_nf(e_n)=f(x).
\end{aligned}
\]
因为 $x$ 是 $V$ 中的任一向量，故
\[
  f=g=b_1f_1+b_2f_2+\cdots+b_nf_n.
\]
这说明 $V^*$ 中任一元素均可表示为 $f_1,f_2,\cdots,f_n$ 的线性组合. 另一方面，若存在 $c_1,c_2,\cdots,c_n$，使
\[
  c_1f_1+c_2f_2+\cdots+c_nf_n=0,
\]
则将上式依次作用于 $e_1,e_2,\cdots,e_n$ 上即得
\[
  c_1=c_2=\cdots=c_n=0.
\]
因此 $\{f_1,f_2,\cdots,f_n\}$ 是 $V^*$ 中线性无关的向量组，组成了 $V^*$ 的一组基，称为 $\{e_1,e_2,\cdots,e_n\}$ 的对偶基. 从这里我们顺便得到结论：
\begin{equation}
  \dim V^*=\dim V=n.
  \tag{10.1.1}
\end{equation}
需要注意的是当 $V$ 是无限维空间时，(10.1.1) 式不再成立.

现在引进一个记号 $\langle,\rangle$：
\[
  \langle f,x\rangle=f(x),
\]
其中 $x\in V,\ f\in V^*$. 容易证明 $\langle,\rangle$ 有下列性质.

\noindent\textbf{记号 $\langle,\rangle$ 的性质}
\[
\begin{gathered}
  (1)\quad \text{若 }\langle f,x\rangle=0\text{ 对一切 }x\in V\text{ 成立，则 }f=0;\\
  (2)\quad \text{若 }\langle f,x\rangle=0\text{ 对一切 }f\in V^*\text{ 成立，则 }x=0.
\end{gathered}
\]

\noindent\textbf{证明}\quad (1) 是显然的. 现来证明 (2). 若 $x\ne0$，把 $x$ 作为 $V$ 的一个基向量并扩充它为 $V$ 的一组基 $\{e_1=x,e_2,\cdots,e_n\}$. 令 $f(e_1)=1,\ f(e_i)=0(i>1)$，则 $f$ 可定义成为 $V$ 上的线性函数，这时 $\langle f,x\rangle\ne0$. $\square$

固定 $f$，则 $\langle f,-\rangle$ 便是 $V$ 上的线性函数且 $\langle f,-\rangle=f$. 同样地，固定 $x$，$\langle-,x\rangle$ 可看成是 $V^*$ 上的线性函数. 事实上，对任意的 $f,g\in V^*$ 及任意的 $k\in\mathbb K$，有
\[
\begin{aligned}
  \langle f+g,x\rangle&=(f+g)(x)=f(x)+g(x)\\
  &=\langle f,x\rangle+\langle g,x\rangle,\\
  \langle kf,x\rangle&=(kf)(x)=kf(x)=k\langle f,x\rangle,
\end{aligned}
\]
因此 $\langle-,x\rangle$ 是 $V^*$ 上的线性函数. 记 $V^*$ 上的全体线性函数为 $(V^*)^*=V^{**}$，则 $V^{**}$ 就是 $V^*$ 的共轭空间. 上面的讨论表明：
\[
  \langle-,x\rangle\in V^{**}.
\]
定义 $V\to V^{**}$ 的映射 $\eta$：
\[
  \eta(x)=\langle-,x\rangle,
\]
则
\[
  \eta(x+y)=\langle-,x+y\rangle=\langle-,x\rangle+\langle-,y\rangle=\eta(x)+\eta(y),
\]
\[
  \eta(kx)=\langle-,kx\rangle=k\langle-,x\rangle=k\eta(x),
\]
因此 $\eta$ 是线性空间 $V\to V^{**}$ 的线性映射. 又若 $\eta(x)=0$（这里的 0 表示 $V^{**}$ 的零向量），即对任意的 $f\in V^*$，$\langle f,x\rangle=0$. 由前面的分析得知，必有 $x=0$，这就是说 $\operatorname{Ker}\eta=0$，因此 $\eta$ 是一个单映射. 这时若 $\dim V=n$，则 $\dim V^{**}=\dim V^*=n$，故 $\eta$ 也是满映射，即 $\eta$ 是线性同构. 如果把 $V$ 与 $V^{**}$ 在这个同构下等同起来，则 $V$ 可以看成是 $V^*$ 的对偶空间. 这样 $V$ 与 $V^*$ 具有平等的地位，它们互为对偶.

现在我们来研究如何把线性映射诱导到对偶空间上去. 设 $\varphi$ 是 $\mathbb K$ 上线性空间 $V$ 到 $\mathbb K$ 上线性空间 $U$ 中的线性映射，$V^*,U^*$ 分别是 $V,U$ 的共轭空间. 设 $g\in U^*$，则 $g\varphi$ 是 $V$ 上的线性函数. 定义 $U^*\to V^*$ 的映射 $\varphi^*$ 如下：
\[
  \varphi^*(g)=g\varphi,
\]
则
\[
  \varphi^*(g_1+g_2)=(g_1+g_2)\varphi=g_1\varphi+g_2\varphi=\varphi^*(g_1)+\varphi^*(g_2),
\]
\[
  \varphi^*(kg)=(kg)\varphi=k(g\varphi)=k\varphi^*(g),
\]
因此 $\varphi^*$ 是线性映射. 我们称 $\varphi^*$ 是线性映射 $\varphi$ 的对偶映射.

\noindent\textbf{定理 10.1.1}\quad 设 $V,U$ 是数域 $\mathbb K$ 上的线性空间，$\varphi$ 是 $V\to U$ 的线性映射，$\varphi^*:U^*\to V^*$ 是 $\varphi$ 的对偶映射，则

(1) 对任意的 $x\in V$ 及任意的 $g\in U^*$，总成立：
\begin{equation}
  \langle\varphi^*(g),x\rangle=\langle g,\varphi(x)\rangle,
  \tag{10.1.2}
\end{equation}
若 $\widetilde\varphi$ 是 $U^*\to V^*$ 的线性映射且等式
\begin{equation}
  \langle\widetilde\varphi(g),x\rangle=\langle g,\varphi(x)\rangle
  \tag{10.1.3}
\end{equation}
对一切 $x\in V,\ g\in U^*$ 成立，那么 $\widetilde\varphi=\varphi^*$.

(2) 若 $V,U$ 是有限维线性空间，设 $\{e_1,e_2,\cdots,e_n\}$ 是 $V$ 的一组基，$\{f_1,f_2,\cdots,f_n\}$ 是其对偶基；$\{u_1,u_2,\cdots,u_m\}$ 是 $U$ 的一组基，$\{g_1,g_2,\cdots,g_m\}$ 是其对偶基. 设 $\varphi$ 在基 $\{e_1,e_2,\cdots,e_n\}$ 和基 $\{u_1,u_2,\cdots,u_m\}$ 下的表示矩阵是 $A$，则 $\varphi^*$ 在基 $\{g_1,g_2,\cdots,g_m\}$ 和基 $\{f_1,f_2,\cdots,f_n\}$ 下的表示矩阵为 $A'$.

\noindent\textbf{证明}\quad (1) 由 $\varphi^*$ 之定义，有
\[
\begin{aligned}
  \langle\varphi^*(g),x\rangle&=\varphi^*(g)(x)=(g\varphi)(x)\\
  &=g(\varphi(x))=\langle g,\varphi(x)\rangle.
\end{aligned}
\]
又若 $\widetilde\varphi$ 使 (10.1.3) 式成立，则对任意的 $x\in V$，有
\[
  \widetilde\varphi(g)(x)=\langle\widetilde\varphi(g),x\rangle=\langle g,\varphi(x)\rangle
  =\langle\varphi^*(g),x\rangle=\varphi^*(g)(x),
\]
因此 $\widetilde\varphi(g)=\varphi^*(g)$ 对一切 $g\in U^*$ 成立，即有 $\widetilde\varphi=\varphi^*$.

(2) 设 $A=(a_{ij})_{m\times n}$，且 $B=(b_{ji})_{n\times m}$ 为 $\varphi^*$ 在对偶基下的表示矩阵，即有
\[
  \varphi(e_j)=\sum_{i=1}^m a_{ij}u_i,\quad j=1,2,\cdots,n,
\]
\[
  \varphi^*(g_i)=\sum_{j=1}^n b_{ji}f_j,\quad i=1,2,\cdots,m.
\]
对任意的 $1\le k\le m,\ 1\le l\le n$，在 (10.1.2) 式中令 $g=g_k,\ x=e_l$，计算可得
\[
  \left\langle\varphi^*(g_k),e_l\right\rangle
  =\left\langle\sum_{j=1}^n b_{jk}f_j,e_l\right\rangle=b_{lk},
\]
\[
  \left\langle g_k,\varphi(e_l)\right\rangle
  =\left\langle g_k,\sum_{i=1}^m a_{il}u_i\right\rangle=a_{kl}.
\]
因此 $b_{lk}=a_{kl}$ 对任意的 $1\le k\le m,\ 1\le l\le n$ 成立，即 $B=A'$. $\square$

\noindent\textbf{推论 10.1.1}\quad 设 $V,U,W$ 是数域 $\mathbb K$ 上的线性空间，$\varphi,\varphi_1,\varphi_2$ 是 $V\to U$ 的线性映射，$\psi$ 是 $U\to W$ 的线性映射，则

(1) $(k_1\varphi_1+k_2\varphi_2)^*=k_1\varphi_1^*+k_2\varphi_2^*$，其中 $k_1,k_2\in\mathbb K$.

(2) $(\psi\varphi)^*=\varphi^*\psi^*$.

(3) 若 $\varphi:V\to U$ 是线性同构，则 $\varphi^*:U^*\to V^*$ 也是线性同构，此时 $(\varphi^*)^{-1}=(\varphi^{-1})^*$.

(4) 若 $V,U$ 都是有限维线性空间，则 $\varphi$ 是单映射的充分必要条件是 $\varphi^*$ 为满映射，$\varphi$ 是满映射的充分必要条件是 $\varphi^*$ 为单映射. 特别，$\varphi$ 是线性同构的充分必要条件是 $\varphi^*$ 也是线性同构.

\noindent\textbf{证明}\quad (1) 和 (2) 都可以用定理 10.1.1 (1) 来证明. 我们只证明 (2)，而把 (1) 留给读者自己证明. 对任意的 $x\in V,\ h\in W^*$，有
\[
\begin{aligned}
  \left\langle(\psi\varphi)^*(h),x\right\rangle
  &=\langle h,\psi\varphi(x)\rangle=\left\langle\psi^*(h),\varphi(x)\right\rangle\\
  &=\left\langle\varphi^*(\psi^*(h)),x\right\rangle
  =\left\langle(\varphi^*\psi^*)(h),x\right\rangle.
\end{aligned}
\]
由定理 10.1.1 (1) 即得 $(\psi\varphi)^*=\varphi^*\psi^*$.

(3) 显然 $(I_V)^*=I_{V^*}$ 成立. 若 $\varphi$ 是线性同构，则 $\varphi^{-1}\varphi=I_V$，由 (2) 可得
\[
  I_{V^*}=(\varphi^{-1}\varphi)^*=\varphi^*(\varphi^{-1})^*.
\]
同理可证 $I_{U^*}=(\varphi^{-1})^*\varphi^*$，故 (3) 的结论成立.

(4) 由推论 4.4.3 知，$\varphi$ 是单映射的充分必要条件是 $\varphi$ 的表示矩阵 $A$ 是列满秩阵，$\varphi$ 是满映射的充分必要条件是 $\varphi$ 的表示矩阵 $A$ 是行满秩阵. 又 $A$ 是列满秩阵（行满秩阵）当且仅当 $A'$ 是行满秩阵（列满秩阵），故由定理 10.1.1 (2) 知 (4) 的结论成立. $\square$

\noindent\textbf{例 10.1.1}\quad 设 $V$ 是数域 $\mathbb K$ 上的 $n$ 维列向量空间，$V'$ 是 $\mathbb K$ 上的 $n$ 维行向量空间. 现利用 $V'$ 定义 $V$ 上的线性函数如下：设 $v\in V,\ u$ 为 $V'$ 中固定的向量，若
\[
  u=(u_1,u_2,\cdots,u_n),\quad v=(v_1,v_2,\cdots,v_n)',
\]
定义
\[
  \langle u,v\rangle=u_1v_1+u_2v_2+\cdots+u_nv_n,
\]
则 $\langle u,-\rangle$ 是 $V$ 上的线性函数（读者不妨验证之）. 另一方面，设 $f$ 是 $V$ 上的任一线性函数，$\{e_1,e_2,\cdots,e_n\}$ 是 $V$ 的标准基，若 $f(e_i)=a_i$，令
\[
  u=(a_1,a_2,\cdots,a_n),
\]
则不难验证对任意的 $v\in V,\ f(v)=\langle u,v\rangle$. 这表明我们可以把 $V'$ 看成是 $V$ 的对偶空间，同理可把 $V$ 看成是 $V'$ 的对偶空间. 又因为
\[
  \langle e_i',e_j\rangle=\delta_{ij},
\]
所以 $\{e_1',e_2',\cdots,e_n'\}$ 是 $\{e_1,e_2,\cdots,e_n\}$ 的对偶基. 设 $A$ 是 $n$ 阶方阵，则 $A$ 定义了 $V$ 上的线性变换 $\varphi$：
\[
  \varphi(v)=Av,
\]
$A$ 也定义了 $V'$ 上的线性变换 $\psi$：
\[
  \psi(u)=uA.
\]
由于
\[
  \langle uA,v\rangle=(uA)v=u(Av)=\langle u,Av\rangle,
\]
故 $\psi=\varphi^*$.

\section*{习题 10.1}

1. 设 $n$ 维线性空间 $V$ 中的一组基 $\{e_1,e_2,\cdots,e_n\}$ 到另一组基 $\{v_1,v_2,\cdots,v_n\}$ 的过渡矩阵为 $P$，求这两组基的对偶基之间的过渡矩阵.

2. 设 $V_1$ 是线性空间 $V$ 的子空间，记
\[
  V_1^\perp=\{f\in V^*\mid \langle f,V_1\rangle=0\}.
\]
求证：$V_1^\perp$ 是 $V^*$ 的子空间，且若 $V_2$ 是 $V$ 的另一子空间，则
\[
  V_1^\perp\cap V_2^\perp=(V_1+V_2)^\perp.
\]

3. 设 $V$ 是 $\mathbb K$ 上的有限维线性空间，$V_1$ 是 $V$ 的子空间，求证：
\[
  \dim V=\dim V_1+\dim V_1^\perp.
\]

4. 设 $V$ 是 $\mathbb K$ 上的有限维线性空间，$V_1,V_2$ 是 $V$ 的子空间，把 $V$ 看成是 $V^*$ 的对偶空间（按本节中的方式），求证：
\[
  (V_1^\perp)^\perp=V_1,\quad (V_1\cap V_2)^\perp=V_1^\perp+V_2^\perp.
\]

5. 设 $V,U$ 是 $\mathbb K$ 上的有限维线性空间，$\varphi,\psi$ 是 $V\to U$ 的线性映射，求证：$\varphi=\psi$ 当且仅当 $\varphi^*=\psi^*$. 利用这一结论以及复习题四的第 3 题给出推论 10.1.1 (4) 的另一证明.

6. 设 $V,U$ 是 $\mathbb K$ 上的有限维线性空间，$\varphi$ 是 $V\to U$ 的线性映射. 求证：若将 $V$ 与 $V^*$，$U$ 与 $U^*$ 看成是互为对偶的空间，则 $(\varphi^*)^*=\varphi$.

\section*{§10.2\quad 双线性型}

在上一节中，我们在两个线性空间 $V^*,V$ 上定义了一个函数 $\langle,\rangle$，它可以看成是一个``双变元''函数，其中一个变元在 $V^*$ 中，另一个变元在 $V$ 中. 这个双变元函数对每个变元都是线性的，即
\[
\begin{gathered}
  \langle f,x+y\rangle=\langle f,x\rangle+\langle f,y\rangle,\\
  \langle f,ky\rangle=k\langle f,y\rangle,\\
  \langle f+g,x\rangle=\langle f,x\rangle+\langle g,x\rangle,\\
  \langle kf,x\rangle=k\langle f,x\rangle.
\end{gathered}
\]
这样的性质称为双线性，$\langle,\rangle$ 称为 $V^*$ 与 $V$ 上的双线性函数. 现在我们要在数域 $\mathbb K$ 上的任意两个线性空间上定义双线性函数.

\noindent\textbf{定义 10.2.1}\quad 设 $U,V$ 是数域 $\mathbb K$ 上的线性空间，$U\times V$ 是它们的积集合. 若存在集合 $U\times V\to\mathbb K$ 的映射 $g$，满足下列条件：

(1) 对任意的 $x,y\in U,\ z\in V,\ k\in\mathbb K$，
\[
  g(x+y,z)=g(x,z)+g(y,z),
\]
\[
  g(kx,z)=kg(x,z);
\]

(2) 对任意的 $x\in U,\ z,w\in V,\ k\in\mathbb K$，
\[
  g(x,z+w)=g(x,z)+g(x,w),
\]
\[
  g(x,kz)=kg(x,z),
\]
则称 $g$ 是 $U$ 与 $V$ 上的双线性函数或双线性型.

显然，若固定 $z\in V$，$g(-,z)$ 是 $U$ 上的线性函数. 同样，若固定 $x\in U$，$g(x,-)$ 是 $V$ 上的线性函数. 显而易见，§10.1 中的 $\langle,\rangle$ 是定义在 $V^*$ 与 $V$ 上的双线性函数.

现在我们要研究双线性型的表示. 因为双线性型对每个变元都是线性的，不难看出双线性函数的值完全被该函数在基向量上的值唯一确定. 设 $U$ 及 $V$ 分别是数域 $\mathbb K$ 上的 $m$ 维与 $n$ 维线性空间，$\{e_1,e_2,\cdots,e_m\}$ 与 $\{v_1,v_2,\cdots,v_n\}$ 分别是 $U$ 与 $V$ 的基，$g$ 是定义在 $U$ 和 $V$ 上的双线性型. 设 $A$ 是一个 $m\times n$ 矩阵，它的第 $(i,j)$ 元素等于 $g(e_i,v_j)$，即
\[
  A=\left(\begin{array}{cccc}
  g(e_1,v_1)&g(e_1,v_2)&\cdots&g(e_1,v_n)\\
  g(e_2,v_1)&g(e_2,v_2)&\cdots&g(e_2,v_n)\\
  \vdots&\vdots&&\vdots\\
  g(e_m,v_1)&g(e_m,v_2)&\cdots&g(e_m,v_n)
  \end{array}\right).
\]
若 $x\in U,\ y\in V$，
\[
  x=\sum_{i=1}^m a_ie_i,\quad y=\sum_{j=1}^n b_jv_j,
\]
则
\begin{equation}
\begin{aligned}
  g(x,y)&=g\left(\sum_{i=1}^m a_ie_i,\sum_{j=1}^n b_jv_j\right)\\
  &=\sum_{i=1}^m\sum_{j=1}^n a_ib_jg(e_i,v_j)\\
  &=(a_1,a_2,\cdots,a_m)A(b_1,b_2,\cdots,b_n)'.
\end{aligned}
\tag{10.2.1}
\end{equation}
(10.2.1) 式告诉我们，若记 $x$ 的坐标向量（用列向量表示）为 $\alpha$，记 $y$ 的坐标向量为 $\beta$，则
\begin{equation}
  g(x,y)=\alpha' A\beta.
  \tag{10.2.2}
\end{equation}
(10.2.1) 式还表明，取定基后，$U$ 和 $V$ 上的双线性型 $g$ 可以决定 $\mathbb K$ 上的 $m\times n$ 矩阵 $A=(g(e_i,v_j))$. 反过来，若取定 $U$ 和 $V$ 的基以及 $\mathbb K$ 上的 $m\times n$ 矩阵 $A$，则由 (10.2.1) 式可以定义 $U$ 和 $V$ 上的双线性函数 $g$. 因此在固定基下，$U$ 和 $V$ 上的双线性函数全体与 $\mathbb K$ 上的 $m\times n$ 矩阵之间有一个一一对应. 它使我们能够用矩阵作为工具来研究双线性型.

我们自然关心这样一件事：若 $U$ 与 $V$ 的基发生变化，同一个双线性函数在不同基下的表示矩阵有什么关系？现设 $\{e_1,e_2,\cdots,e_m\}$，$\{e_1',e_2',\cdots,e_m'\}$ 是 $U$ 的两组基，$\{v_1,v_2,\cdots,v_n\}$，$\{v_1',v_2',\cdots,v_n'\}$ 是 $V$ 的两组基，它们之间的过渡矩阵分别为 $C$ 及 $D$. 又设 $x\in U$ 在基 $\{e_i\}$ 下的坐标向量为 $\alpha$，在基 $\{e_i'\}$ 下的坐标向量为 $\beta$，则 $\alpha=C\beta$. 类似地，设 $y\in V$ 在基 $\{v_j\}$ 下的坐标向量为 $\gamma$，在基 $\{v_j'\}$ 下的坐标向量为 $\delta$，则 $\gamma=D\delta$. 设双线性型 $g$ 在基 $\{e_i\}$ 和基 $\{v_j\}$ 下的表示矩阵为 $A$，在基 $\{e_i'\}$ 和基 $\{v_j'\}$ 下的表示矩阵为 $B$，于是
\[
  g(x,y)=\alpha' A\gamma=(C\beta)'A(D\delta)=\beta'(C'AD)\delta.
\]
另一方面
\[
  g(x,y)=\beta'B\delta.
\]
因为 $x,y$ 是任意的，所以
\[
  B=C'AD.
\]
这说明，$g$ 在不同基下的表示矩阵是相抵的. 由矩阵理论知道存在非异阵 $C$ 及 $D$，使 $C'AD$ 为相抵标准型. 这就是说，我们可以选择 $U,V$ 的基，使 $g$ 在这两组基下的表示矩阵是相抵标准型
\[
  \left(\begin{array}{cc}
  I_r&O\\
  O&O
  \end{array}\right).
\]
矩阵 $A$ 的秩 $r$ 称为 $g$ 的秩，它不随基的改变而改变. 这样就证明了下面的定理.

\noindent\textbf{定理 10.2.1}\quad 设 $g$ 是 $U$ 和 $V$ 上的双线性型，则必存在 $U$ 的基 $\{e_1,e_2,\cdots,e_m\}$ 及 $V$ 的基 $\{v_1,v_2,\cdots,v_n\}$，使
\[
  g(e_i,v_j)=\delta_{ij},\quad i,j=1,\cdots,r;
\]
\[
  g(e_i,v_j)=0,\quad i>r\text{ 或 }j>r,
\]
其中 $r$ 为 $g$ 的秩.

\noindent\textbf{定义 10.2.2}\quad 设 $g$ 是 $U$ 和 $V$ 上的双线性型. 令
\[
  L=\{u\in U\mid g(u,y)=0\text{ 对一切 }y\in V\text{ 成立}\},
\]
\[
  R=\{v\in V\mid g(x,v)=0\text{ 对一切 }x\in U\text{ 成立}\},
\]
则 $L$ 与 $R$ 分别是 $U$ 与 $V$ 的子空间，分别称为 $g$ 的左根子空间与右根子空间.

设 $U$ 是 $m$ 维线性空间，$V$ 是 $n$ 维线性空间，且 $g$ 的秩等于 $r$，则由定理 10.2.1 知道可选择 $U$ 和 $V$ 的基 $\{e_i\}_{i=1,\cdots,m}$ 及 $\{v_j\}_{j=1,\cdots,n}$，使左根子空间 $L$ 恰好由 $e_{r+1},\cdots,e_m$ 张成，右根子空间 $R$ 恰好由 $v_{r+1},\cdots,v_n$ 张成. 因此我们有
\begin{equation}
  \dim L=m-r,\quad \dim R=n-r.
  \tag{10.2.3}
\end{equation}

\noindent\textbf{定义 10.2.3}\quad 若双线性型 $g$ 的左、右根子空间都为零，则 $g$ 称为非退化的.

\noindent\textbf{定理 10.2.2}\quad $U$ 和 $V$ 上的双线性型 $g$ 为非退化双线性型的充分必要条件是
\[
  \dim U=\dim V=r,
\]
其中 $r$ 为 $g$ 的秩.

\noindent\textbf{证明}\quad 由 (10.2.3) 式即得. $\square$

\noindent\textbf{推论 10.2.1}\quad $U$ 和 $V$ 上的双线性型 $g$ 为非退化的充分必要条件是 $g$ 在 $U$ 和 $V$ 的任意两组基下的表示矩阵都是非异阵.

现设 $g$ 是 $U$ 和 $V$ 上非退化的双线性型，固定 $x$，$g(x,-)$ 就成了 $V$ 上的线性函数. 作 $U\to V^*$ 的映射 $\varphi$：
\[
  \varphi(x)=g(x,-),
\]
则 $\varphi$ 是一个线性映射. 若 $\varphi(x)=0$，即 $g(x,y)=0$ 对一切 $y\in V$ 成立，由于 $g$ 是非退化的，故 $x=0$. 这说明 $\operatorname{Ker}\varphi=0$，也就是 $\varphi$ 为 $U\to V^*$ 的单映射. 另一方面，$\dim V^*=\dim V=\dim U$，故 $\varphi$ 是一个线性同构. 若将 $x$ 与 $g(x,-)$ 等同起来，则 $U$ 就成了 $V$ 的对偶空间 $V^*$，这时有
\[
  \langle\varphi(x),y\rangle=g(x,y).
\]
类似地，我们可将 $V$ 与 $U^*$ 等同起来，即存在线性同构 $\psi:V\to U^*$，使
\[
  \langle x,\psi(y)\rangle=g(x,y)
\]
对一切 $x\in U,\ y\in V$ 成立. 这一事实表明非退化双线性型在同构的意义下与 §10.1 中定义的双线性型 $\langle,\rangle$ 没有什么区别，许多问题的讨论可归结为 $V^*$ 和 $V$ 上的双线性型 $\langle,\rangle$ 的讨论.

一般地，我们还有下列定理.

\noindent\textbf{定理 10.2.3}\quad 设 $g_1$ 及 $g_2$ 是 $U$ 和 $V$ 上的两个非退化双线性型，则存在 $U$ 上的非异线性变换 $\varphi$ 及 $V$ 上的非异线性变换 $\psi$，使
\[
  g_2(\varphi(x),y)=g_1(x,y),\quad g_2(x,\psi(y))=g_1(x,y)
\]
对一切 $x\in U,\ y\in V$ 成立.

\noindent\textbf{证明}\quad 由上面的讨论知道存在线性同构 $\varphi_1:U\to V^*$，使
\[
  \langle\varphi_1(x),y\rangle=g_1(x,y).
\]
同理，存在线性同构 $\varphi_2:U\to V^*$，使
\[
  \langle\varphi_2(x),y\rangle=g_2(x,y).
\]
令 $\varphi=\varphi_2^{-1}\varphi_1$，则 $\varphi$ 是 $U$ 上的非异线性变换，且
\[
\begin{aligned}
  g_2(\varphi(x),y)&=\langle\varphi_2\varphi(x),y\rangle\\
  &=\langle\varphi_1(x),y\rangle\\
  &=g_1(x,y).
\end{aligned}
\]
同理可证明另一等式. $\square$

对偶映射的概念也可推广到双线性型上. 设 $g_1$ 及 $g_2$ 是 $U$ 和 $V$ 上的非退化双线性型，$\varphi$ 是 $V$ 上的线性变换，若存在 $U$ 上的线性变换 $\varphi^*$，使
\[
  g_2(\varphi^*(x),y)=g_1(x,\varphi(y))
\]
对一切 $x\in U,\ y\in V$ 成立，则称 $\varphi^*$ 是 $\varphi$ 关于 $g_1,g_2$ 的对偶. 不难证明 $\varphi$ 关于 $g_1,g_2$ 对偶的存在性. 事实上，令 $\varphi^*=\varphi_2^{-1}\widetilde\varphi\varphi_1$，其中 $\varphi_1$ 和 $\varphi_2$ 是定理 10.2.3 证明中的同构，$\widetilde\varphi$ 是 $\varphi$ 关于 $\langle,\rangle$ 的对偶，即 $\widetilde\varphi$ 适合
\[
  \langle\widetilde\varphi(v),y\rangle=\langle v,\varphi(y)\rangle,
\]
于是
\[
\begin{aligned}
  g_2(\varphi^*(x),y)&=\langle\varphi_2\varphi^*(x),y\rangle
  =\langle\widetilde\varphi\varphi_1(x),y\rangle\\
  &=\langle\varphi_1(x),\varphi(y)\rangle
  =g_1(x,\varphi(y)).
\end{aligned}
\]
唯一性可这样来证明：若 $\xi$ 是 $U$ 上的线性变换，且
\[
  g_2(\xi(x),y)=g_1(x,\varphi(y)),
\]
则
\[
  g_2(\xi(x),y)=g_2(\varphi^*(x),y)
\]
对一切 $x\in U,\ y\in V$ 成立，因此
\[
  g_2(\xi(x)-\varphi^*(x),y)=0
\]
对一切 $y\in V$ 成立. 由于 $g_2$ 是非退化的双线性型，故 $\xi(x)-\varphi^*(x)=0$，即 $\xi(x)=\varphi^*(x)$ 对一切 $x\in U$ 成立，因此 $\xi=\varphi^*$.

\section*{习题 10.2}

1. 设 $g_1$ 及 $g_2$ 是 $\mathbb K$ 上的线性空间 $U$ 和 $V$ 上的双线性型，定义
\[
  (g_1+g_2)(x,y)=g_1(x,y)+g_2(x,y),
\]
对 $\mathbb K$ 中任一数 $k$，定义
\[
  (kg_1)(x,y)=kg_1(x,y).
\]
求证：$g_1+g_2$ 和 $kg_1$ 仍是 $U$ 和 $V$ 上的双线性型，且在此定义下，$U$ 和 $V$ 上的全体双线性型成为 $\mathbb K$ 上的线性空间. 若 $\dim U=m,\ \dim V=n$，则上述线性空间的维数为 $mn$.

2. 设 $g$ 是 $U$ 和 $V$ 上的双线性型，$V^*$ 是 $V$ 的对偶空间，$L$ 是 $g$ 的左根子空间. 定义 $U\to V^*$ 的映射 $\varphi$：
\[
  \varphi(x)=g(x,-),
\]
求证：$\varphi$ 是线性映射且 $\operatorname{Ker}\varphi=L$.

3. 设 $g$ 是 $U$ 和 $V$ 上的非退化双线性型. 若 $\{e_1,e_2,\cdots,e_n\}$，$\{v_1,v_2,\cdots,v_n\}$ 分别是 $U,V$ 的一组基，使
\[
  g(e_i,v_j)=\delta_{ij},\quad i,j=1,2,\cdots,n,
\]
则称 $\{e_i\}$，$\{v_j\}$ 是关于 $g$ 的对偶基. 设 $\varphi$ 是 $V$ 上的线性变换，$\varphi^*$ 是 $\varphi$ 关于 $g$ 的对偶映射. 证明：若 $\varphi$ 在 $\{v_i\}$ 下的表示矩阵为 $A$，则 $\varphi^*$ 在 $\{e_i\}$ 下的表示矩阵为 $A'$.

4. 设 $g$ 是 $U$ 和 $V$ 上的非零双线性型，证明：必存在 $U,V$ 的子空间 $U_0,V_0$，使 $g$ 限制在 $U_0,V_0$ 上是非退化的双线性型，且
\[
  \dim U_0=\dim V_0=\dim U-\dim L,
\]
其中 $L$ 是 $g$ 的左根子空间.

5. 设 $g$ 是 $U$ 和 $V$ 上的非退化双线性型，$S,T$ 分别是 $U,V$ 的子集，且
\[
  S^\perp=\{v\in V\mid g(S,v)=0\},\quad T^\perp=\{u\in U\mid g(u,T)=0\},
\]
求证：$S^\perp$ 是 $V$ 的子空间，$T^\perp$ 是 $U$ 的子空间.

6. 试将 §10.1 习题 2～习题 4 推广到 $U$ 和 $V$ 上的非退化双线性型的情形.

\section*{§10.3\quad 纯量积}

\noindent\textbf{定义 10.3.1}\quad 设 $g$ 是 $V\times V\to\mathbb K$ 的双线性函数，则称 $g$ 是 $V$ 上的纯量积或数量积. 若
\[
  g(x,y)=g(y,x)
\]
对一切 $x,y\in V$ 成立，则称 $g$ 是 $V$ 上的对称型. 若
\[
  g(x,y)=-g(y,x)
\]
对一切 $x,y\in V$ 成立，则称 $g$ 是 $V$ 上的交错型.

交错型又称为反对称型. 交错型的另一等价说法是，对 $V$ 中任一元素 $x$，$g(x,x)=0$. 事实上，由 $g(x,x)=-g(x,x)$ 可推出 $g(x,x)=0$. 另一方面，若 $g(x,x)=0$ 对一切 $x$ 成立，则 $g(x+y,x+y)=0$，由此即可推出 $g(x,y)=-g(y,x)$.

$V$ 上的纯量积也可用矩阵来表示，但是与一般的双线性型的矩阵表示有一点区别. 我们只取 $V$ 的一组基而不是取两组基. 因此若 $\{e_1,e_2,\cdots,e_n\}$ 是 $V$ 的一组基，则 $g$ 的表示矩阵为
\[
  A=(g(e_i,e_j)).
\]
若
\[
  x=\sum_{i=1}^n a_ie_i,\quad y=\sum_{i=1}^n b_ie_i,
\]
则
\begin{equation}
  g(x,y)=(a_1,a_2,\cdots,a_n)A(b_1,b_2,\cdots,b_n)'.
  \tag{10.3.1}
\end{equation}
设 $\{v_i\}$ 是 $V$ 的另一组基，$B$ 是 $g$ 在 $\{v_i\}$ 下的表示矩阵，若从 $\{e_i\}$ 到 $\{v_i\}$ 的过渡矩阵为 $C$，则
\[
  B=C'AC.
\]
这就是说纯量积 $g$ 在不同基下的表示矩阵是合同的.

显然，对称型的表示矩阵是对称阵，反对称型的表示矩阵是反对称阵. 反过来，在取定 $V$ 的一组基后，若给定一个对称阵（反对称阵），则 (10.3.1) 式定义了 $V$ 上的一个对称型（反对称型）.

设 $g$ 是 $V$ 上的纯量积，$x,y$ 是 $V$ 中两个向量. 若 $g(x,y)=0$，则称 $x$ 左垂直于 $y$ 或 $y$ 右垂直于 $x$，记为 $x\perp y$. 当 $g$ 是对称型或交错型时，$x\perp y$ 等价于 $y\perp x$，这时称 $x$ 与 $y$ 正交. 反过来，如果 $g$ 有下列性质：由 $x\perp y$ 总可推出 $y\perp x$，那么 $g$ 是否必为对称型或交错型呢？下面的定理将给予肯定的回答.

\noindent\textbf{定理 10.3.1}\quad 设 $g$ 是 $V$ 上的纯量积，则在 $V$ 中 $x\perp y$ 等价于 $y\perp x$ 的充分必要条件是 $g$ 为对称型或交错型.

\noindent\textbf{证明}\quad 只需证明必要性. 设 $x,y,z\in V$，令
\[
  w=g(x,y)z-g(x,z)y,
\]
则
\[
\begin{aligned}
  g(x,w)&=g(x,g(x,y)z-g(x,z)y)\\
  &=g(x,y)g(x,z)-g(x,z)g(x,y)\\
  &=0,
\end{aligned}
\]
\[
\begin{aligned}
  g(w,x)&=g(g(x,y)z-g(x,z)y,x)\\
  &=g(x,y)g(z,x)-g(x,z)g(y,x).
\end{aligned}
\]
由于从 $x\perp w$ 可推出 $w\perp x$，故
\begin{equation}
  g(x,y)g(z,x)=g(y,x)g(x,z)
  \tag{10.3.2}
\end{equation}
对任意的 $x,y,z\in V$ 均成立. 令 $x=y$ 得
\begin{equation}
  g(x,x)(g(x,z)-g(z,x))=0.
  \tag{10.3.3}
\end{equation}
因此或者 $g(x,x)=0$ 或者 $g(x,z)=g(z,x)$. 若 $g(x,x)=0$ 对一切 $x\in V$ 成立，则 $g$ 是交错型. 我们的目的是要证明 (10.3.3) 式或者对所有的 $x$ 有 $g(x,x)=0$，或者对所有的 $x,z$ 有 $g(x,z)=g(z,x)$. 假设不然，则存在 $u,v,w\in V$，使 $g(u,v)\ne g(v,u),\ g(w,w)\ne0$. 由 (10.3.3) 式知道，$g(u,u)=g(v,v)=0$，并且
\[
  g(w,u)=g(u,w),\quad g(w,v)=g(v,w).
\]
因为 $g(u,v)\ne g(v,u)$，由 (10.3.2) 式知道 $g(u,w)=g(w,u)=0$ 以及 $g(v,w)=g(w,v)=0$，于是
\[
  g(u,w+v)=g(u,v)\ne g(v,u)=g(w+v,u).
\]
由 (10.3.3) 式得 $g(w+v,w+v)=0$，但是
\[
  g(w+v,w+v)=g(w,w)+g(v,v)+g(w,v)+g(v,w)=g(w,w)\ne0,
\]
这就出现了矛盾. $\square$

\noindent\textbf{推论 10.3.1}\quad 若 $g$ 是 $V$ 上的对称型或交错型，$U$ 是 $V$ 的子空间，记
\[
  U^{\perp l}=\{x\in V\mid g(x,U)=0\},
\]
\[
  U^{\perp r}=\{x\in V\mid g(U,x)=0\},
\]
则 $U^{\perp l}=U^{\perp r}$. 特别地，$g$ 的左根子空间等于右根子空间.

\noindent\textbf{注}\quad 这时记 $U^\perp=U^{\perp l}=U^{\perp r}$，称为 $U$ 的正交补空间. $g$ 的左、右根子空间重合，称为 $g$ 的根子空间.

需要注意的是，这里的正交补概念比内积空间中的正交补概念更一般. 因为实线性空间 $V$ 上的内积显然是一种纯量积，而内积空间中正交向量的有关性质在这里不一定成立. 比如在内积空间中，子空间与其正交补之交必是零空间. 但即使在非退化纯量积的情形，子空间 $U$ 与其正交补 $U^\perp$ 之交未必为零.

例如，设 $V$ 是二维实线性空间，$e_1,e_2$ 是 $V$ 的一组基，$g$ 为由矩阵
\[
  \left(\begin{array}{cc}
  1&0\\
  0&-1
  \end{array}\right)
\]
定义的纯量积，显然 $g$ 是非退化的对称型. 令 $u=e_1+e_2$，则 $g(u,u)=0,\ u\perp u$！若记 $U$ 为 $u$ 生成的子空间，则有 $U\perp U$.

\noindent\textbf{定理 10.3.2}\quad 设 $g$ 是 $n$ 维线性空间 $V$ 上的对称型（交错型），$U$ 是 $V$ 的子空间，则 $U\cap U^\perp=0$ 的充分必要条件是 $g$ 限制在 $U$ 上是 $U$ 的一个非退化的纯量积. 这时有直和分解：
\[
  V=U\oplus U^\perp.
\]

\noindent\textbf{证明}\quad 设 $U\cap U^\perp=0$，若 $g$ 限制在 $U$ 上退化，则存在 $U$ 中非零向量 $u$，使得 $g(u,U)=0$，从而 $u\in U\cap U^\perp$，推出矛盾. 反之，设 $g$ 限制在 $U$ 上非退化，任取 $u\in U\cap U^\perp$，则 $g(u,U)=0$，从而 $u=0$，这表明 $U\cap U^\perp=0$.

对于第二个结论，我们先证明若 $g$ 限制在 $U$ 上非退化，则
\[
  \dim U+\dim U^\perp=n.
\]
对任意的 $v\in V,\ g(v,-)$ 限制在 $U$ 上是 $U$ 上的线性函数. 作线性映射
\[
  \varphi:V\to U^*,\quad \varphi(v)=g(v,-),
\]
则 $\operatorname{Ker}\varphi=U^\perp$. 因为 $g$ 限制在 $U$ 上非退化，故限制映射 $\varphi|_U:U\to U^*$ 是单映射，又 $\dim U=\dim U^*$，从而 $\varphi|_U:U\to U^*$ 是同构. 因此对任意的 $f\in U^*$，存在 $u\in U$，使 $f=g(u,-)$，于是 $\varphi$ 是满映射. 最后，由线性映射的维数公式即得 $\dim U+\dim U^\perp=n$，又因为 $U\cap U^\perp=0$，所以 $V=U\oplus U^\perp$. $\square$

\noindent\textbf{注}\quad 当 $g$ 非退化时，有 $(U^\perp)^\perp=U$，故由定理 10.3.2 知，此时 $g$ 限制在 $U$ 上非退化当且仅当 $g$ 限制在 $U^\perp$ 上也非退化.

\noindent\textbf{定理 10.3.3}\quad 设 $g$ 与 $h$ 是 $V$ 上的两个非退化的纯量积，则存在 $V$ 上唯一的非异线性变换 $\varphi$，使
\[
  h(x,y)=g(\varphi(x),y)
\]
对一切 $x,y\in V$ 成立.

\noindent\textbf{证明}\quad 我们用矩阵方法来证明这一结论. 取 $V$ 的一组基 $\{e_1,e_2,\cdots,e_n\}$，设
\[
  A=(h(e_i,e_j)),\quad B=(g(e_i,e_j)),
\]
\[
  x=\sum_{i=1}^n a_ie_i,\quad y=\sum_{i=1}^n b_ie_i,
\]
则
\[
  h(x,y)=(a_1,a_2,\cdots,a_n)A(b_1,b_2,\cdots,b_n)',
\]
\[
  g(x,y)=(a_1,a_2,\cdots,a_n)B(b_1,b_2,\cdots,b_n)'.
\]
设 $\varphi$ 在给定基下的表示矩阵为 $C$，满足 $C'=AB^{-1}$，则 $\varphi(x)$ 的坐标向量为 $C(a_1,a_2,\cdots,a_n)'$，于是
\[
\begin{aligned}
  g(\varphi(x),y)&=(a_1,a_2,\cdots,a_n)C'B(b_1,b_2,\cdots,b_n)'\\
  &=(a_1,a_2,\cdots,a_n)A(b_1,b_2,\cdots,b_n)'\\
  &=h(x,y).
\end{aligned}
\]
唯一性的证明留给读者完成. $\square$

\section*{习题 10.3}

1. 设 $g$ 是 $V$ 上的纯量积（不一定非退化），$\varphi$ 是 $V$ 上的线性变换. 设 $\{e_1,e_2,\cdots,e_n\}$ 是 $V$ 的一组基，$g$ 在这组基下的表示矩阵为 $A=(g(e_i,e_j))$，$\varphi$ 在这组基下的表示矩阵为 $B$. 求证：$g(\varphi(x),y)$ 是 $V$ 上的纯量积且它在这组基下的表示矩阵为 $B'A$.

2. 设 $g$ 与 $h$ 是 $V$ 上秩相同的纯量积，求证：必存在 $V$ 上的非异线性变换 $\varphi$ 和 $\psi$，使
\[
  h(x,y)=g(\varphi(x),\psi(y))
\]
对一切 $x,y\in V$ 成立.

3. 求证：$V$ 上任一纯量积均可表示为一个对称型与一个交错型之和.

4. 设 $W=U\oplus V$，$g$ 及 $h$ 分别是 $U$ 及 $V$ 上的纯量积. 现定义 $W$ 上的纯量积 $q$ 如下：
\[
  q(x+y,u+v)=g(x,u)+h(y,v),
\]
其中 $x,u\in U,\ y,v\in V$. 求证：

(1) 若 $g,h$ 非退化，则 $q$ 也非退化；

(2) 若 $g,h$ 为对称型（交错型），则 $q$ 也为对称型（交错型）；

(3) 若 $\{e_i\}_{i=1,\cdots,m}$，$\{v_i\}_{i=1,\cdots,n}$ 分别是 $U,V$ 的一组基且 $g,h$ 在这两组基下的表示矩阵分别为 $A,B$，则 $q$ 在 $W$ 的基 $\{e_i\}\cup\{v_i\}$ 下的表示矩阵为
\[
  \left(\begin{array}{cc}
  A&O\\
  O&B
  \end{array}\right).
\]

\section*{§10.4\quad 交错型与辛几何}

\noindent\textbf{定理 10.4.1}\quad 设 $g$ 是 $n$ 维线性空间 $V$ 上的交错型，则存在 $V$ 的一组基
\[
  \{u_1,v_1,u_2,v_2,\cdots,u_r,v_r;w_1,\cdots,w_{n-2r}\},
\]
使 $g$ 在这组基下的表示矩阵为分块对角阵：
\begin{equation}
  \operatorname{diag}\{S,S,\cdots,S;0,\cdots,0\},
  \tag{10.4.1}
\end{equation}
其中共有 $r$ 个二阶方阵 $S$，$S$ 为如下形状：
\[
  \left(\begin{array}{cc}
  0&1\\
  -1&0
  \end{array}\right).
\]

\noindent\textbf{证明}\quad 不妨设 $g\ne0$. 这时有 $u,v\in V$，使 $g(u,v)=a\ne0$. 令 $u_1=a^{-1}u,\ v_1=v$，则 $g(u_1,v_1)=-g(v_1,u_1)=1$. 又 $u_1,v_1$ 必线性无关，否则将有 $g(u_1,v_1)=0$. 现设已选定 $u_1,v_1,\cdots,u_j,v_j$ 共 $2j$ 个线性无关的向量，适合
\[
  g(u_i,v_i)=1=-g(v_i,u_i),\quad i=1,\cdots,j;
\]
\[
  g(u_i,y)=0,\quad y\in\{u_1,v_1,\cdots,u_j,v_j\}\setminus\{v_i\};
\]
\[
  g(x,v_i)=0,\quad x\in\{u_1,v_1,\cdots,u_j,v_j\}\setminus\{u_i\}.
\]
设 $V_0$ 为由 $\{u_1,v_1,\cdots,u_j,v_j\}$ 生成的子空间，令
\[
  V_0^\perp=\{v\in V\mid g(v,V_0)=0\},
\]
注意到 $g$ 限制在 $V_0$ 上是非退化的，故由定理 10.3.2 得到 $V=V_0\oplus V_0^\perp$. 若 $g$ 在 $V_0^\perp$ 上的限制为零，则定理已得证. 若 $g$ 在 $V_0^\perp$ 上的限制不等于零，则可找到 $u_{j+1},v_{j+1}\in V_0^\perp$，使
\[
  g(u_{j+1},v_{j+1})=1=-g(v_{j+1},u_{j+1}).
\]
不断这样做下去，便可找到我们所需要的基. $\square$

\noindent\textbf{注}\quad 上述定理也可以用矩阵方法证明，参考 §8.1 习题 8.

\noindent\textbf{推论 10.4.1}\quad 数域 $\mathbb K$ 上的反对称阵的秩必为偶数，且它的行列式等于 $\mathbb K$ 中某个元素的平方.

\noindent\textbf{推论 10.4.2}\quad 数域 $\mathbb K$ 上的两个 $n$ 阶反对称阵合同的充分必要条件是它们具有相同的秩.

\noindent\textbf{定义 10.4.1}\quad 设 $V$ 是 $\mathbb K$ 上的有限维线性空间，若 $V$ 上定义了一个非退化的交错型，则称 $V$ 是一个辛空间.

由推论 10.4.1 知道辛空间的维数总是偶数，且由定理 10.4.1 可知任一 $n$ 维辛空间总有一组基 $\{u_1,v_1,u_2,v_2,\cdots,u_r,v_r\}$，使 $g$ 在这组基下的表示矩阵为分块对角阵：
\[
  \operatorname{diag}\{S,S,\cdots,S\},
\]
其中 $n=2r$，$S$ 为如下形状的二阶方阵：
\[
  \left(\begin{array}{cc}
  0&1\\
  -1&0
  \end{array}\right).
\]
我们称这样的基为 $V$ 的辛基. 在辛基下，$g$ 的形式十分简单. 设
\begin{equation}
  x=a_1u_1+b_1v_1+a_2u_2+b_2v_2+\cdots+a_ru_r+b_rv_r,
  \tag{10.4.2}
\end{equation}
\begin{equation}
  y=a_1'u_1+b_1'v_1+a_2'u_2+b_2'v_2+\cdots+a_r'u_r+b_r'v_r,
  \tag{10.4.3}
\end{equation}
则
\begin{equation}
  g(x,y)=a_1b_1'-a_1'b_1+a_2b_2'-a_2'b_2+\cdots+a_rb_r'-a_r'b_r.
  \tag{10.4.4}
\end{equation}

\noindent\textbf{定义 10.4.2}\quad 设 $V$ 是一个辛空间，$\varphi$ 是 $V$ 上的非异线性变换，若
\[
  g(\varphi(x),\varphi(y))=g(x,y)
\]
对一切 $x,y\in V$ 成立，则 $\varphi$ 称为 $V$ 上的一个辛变换.

辛变换有下列性质.

\noindent\textbf{定理 10.4.2}\quad 设 $V$ 是数域 $\mathbb K$ 上的辛空间，则

(1) $V$ 上的线性变换 $\varphi$ 是辛变换的充分必要条件是 $\varphi$ 将辛基变为辛基；

(2) 两个辛变换之积仍是辛变换；

(3) 恒等变换是辛变换；

(4) 辛变换的逆变换是辛变换.

\noindent\textbf{证明}\quad (1) 必要性显然成立，下面证充分性. 设 $\varphi$ 将辛基变为辛基，则 $\varphi$ 是非异线性变换. 设 $x,y$ 如 (10.4.2) 式及 (10.4.3) 式所示，又 $u_i'=\varphi(u_i),\ v_i'=\varphi(v_i)$，则 $\{u_1',v_1',u_2',v_2',\cdots,u_r',v_r'\}$ 构成 $V$ 的一组辛基，且
\[
  \varphi(x)=a_1u_1'+b_1v_1'+a_2u_2'+b_2v_2'+\cdots+a_ru_r'+b_rv_r',
\]
\[
  \varphi(y)=a_1'u_1'+b_1'v_1'+a_2'u_2'+b_2'v_2'+\cdots+a_r'u_r'+b_r'v_r'.
\]
由 (10.4.4) 式即知
\[
  g(\varphi(x),\varphi(y))=g(x,y).
\]
(2)、(3) 和 (4) 由辛变换的定义直接验证即得，下面以 (4) 为例. 设 $\varphi^{-1}$ 是 $\varphi$ 的逆变换，则
\[
  g(\varphi^{-1}(x),\varphi^{-1}(y))
  =g(\varphi(\varphi^{-1}(x)),\varphi(\varphi^{-1}(y)))=g(x,y),
\]
故 $\varphi^{-1}$ 也是辛变换. $\square$

辛空间上的几何学称为辛几何. 辛几何近年来发展很快，并在数学、力学及其他学科中有重要的应用.

\section*{习题 10.4}

1. 设 4 维辛空间 $(V,g)$ 在一组基 $\{e_1,e_2,e_3,e_4\}$ 下的表示矩阵为
\[
  A=\left(\begin{array}{rrrr}
  0&2&-1&3\\
  -2&0&4&-2\\
  1&-4&0&1\\
  -3&2&-1&0
  \end{array}\right),
\]
求 $V$ 的一组辛基.

2. 设 $V$ 是辛空间，$\varphi$ 是 $V$ 上的辛变换，$\varphi$ 在一组辛基下的表示矩阵称为辛矩阵. 令
\[
  A=\operatorname{diag}\{S,S,\cdots,S\},
\]
其中有 $r$ 个 $S$ 且
\[
  S=\left(\begin{array}{rr}
  0&1\\
  -1&0
  \end{array}\right).
\]
求证：$n(n=2r)$ 阶方阵 $T$ 是辛矩阵的充分必要条件是
\[
  T'AT=A.
\]

3. 证明：辛变换的行列式的值等于 1 或 $-1$.

\section*{§10.5\quad 对称型与正交几何}

设 $V$ 是数域 $\mathbb K$ 上的线性空间，$g$ 是 $V$ 上的对称型. 取 $V$ 的一组基 $\{e_1,e_2,\cdots,e_n\}$，设 $g$ 在这组基下的表示矩阵为 $A$，则 $A$ 是对称阵. 若
\[
  x=x_1e_1+x_2e_2+\cdots+x_ne_n,
\]
则
\[
  g(x,x)=(x_1,x_2,\cdots,x_n)A(x_1,x_2,\cdots,x_n)'.
\]
这时我们得到了一个 $V$ 上的二次型. 反过来，若 $f$ 是 $V$ 上的二次型，定义
\[
  g(x,y)=\frac12\bigl(f(x+y)-f(x)-f(y)\bigr).
\]
在基 $\{e_1,e_2,\cdots,e_n\}$ 下，设二次型 $f$ 的相伴对称阵为 $A$. 若
\[
  x=x_1e_1+x_2e_2+\cdots+x_ne_n,\quad y=y_1e_1+y_2e_2+\cdots+y_ne_n,
\]
$\alpha,\beta$ 分别是 $x,y$ 在给定基下的坐标向量，即
\[
  \alpha=(x_1,x_2,\cdots,x_n)',\quad \beta=(y_1,y_2,\cdots,y_n)',
\]
则
\[
\begin{aligned}
  g(x,y)&=\frac12\left((\alpha'+\beta')A(\alpha+\beta)-\alpha'A\alpha-\beta'A\beta\right)\\
  &=\frac12(\alpha'A\alpha+\beta'A\beta+\beta'A\alpha+\alpha'A\beta-\alpha'A\alpha-\beta'A\beta)\\
  &=\frac12(\alpha'A\beta+\beta'A\alpha).
\end{aligned}
\]
因为 $A$ 是对称阵，故 $\alpha'A\beta=\beta'A\alpha$，于是
\[
  g(x,y)=g(y,x)=\alpha'A\beta.
\]
显然，$g$ 是 $V$ 上的对称型. 从这里可以看出，对称型与二次型是等价的. 在这个意义下，双线性型可以看成是二次型的推广.

我们已经知道，任一 $\mathbb K$ 上的对称阵必合同于 $\mathbb K$ 上的对角阵. 因此，若 $g$ 是 $V$ 上的对称型，则必存在 $V$ 的一组基，使 $g$ 在这组基下的表示矩阵为对角阵
\[
  \operatorname{diag}\{b_1,\cdots,b_r;0,\cdots,0\}\quad (b_i\ne0).
\]
这组基称为 $V$ 的正交基.

欧氏空间的内积实际上就是一个正定对称型. 若一个线性空间 $V$ 上定义了一个非退化的对称型，则 $V$ 上的几何学称为正交几何学. 正交几何是欧氏几何的推广. 下面我们来研究正交几何学中的一些基本结论.

\noindent\textbf{定义 10.5.1}\quad 设 $V_i(i=1,2)$ 是有限维线性空间，$g_i(i=1,2)$ 是 $V_i$ 上非退化的对称型. 若存在 $V_1\to V_2$ 的线性同构 $\eta$，使
\[
  g_2(\eta(x),\eta(y))=g_1(x,y)
\]
对一切 $x,y\in V_1$ 成立，则称 $\eta$ 是 $V_1\to V_2$ 的保距同构. 特别，当 $V_1=V_2=V,\ g_1=g_2=g$ 时，则 $\eta$ 称为 $(V,g)$ 上的一个正交变换.

\noindent\textbf{定理 10.5.1}\quad 设 $V$ 是数域 $\mathbb K$ 上的有限维线性空间，$g$ 是 $V$ 上非退化的对称型，则

(1) $V$ 上两个正交变换之积仍是正交变换；

(2) $V$ 上恒等映射是正交变换；

(3) $V$ 上正交变换的逆变换也是正交变换.

\noindent\textbf{证明}\quad 由正交变换的定义直接验证即得. $\square$

给定 $u\in V$ 且 $g(u,u)\ne0$，构造下列映射 $S_u$：
\[
  S_u(x)=x-\frac{2g(x,u)}{g(u,u)}u,
\]
我们有
\[
\begin{aligned}
  g(S_u(x),S_u(y))
  &=g\left(x-\frac{2g(x,u)}{g(u,u)}u,\ y-\frac{2g(y,u)}{g(u,u)}u\right)\\
  &=g(x,y)-\frac{2g(y,u)}{g(u,u)}g(x,u)\\
  &\quad-\frac{2g(x,u)}{g(u,u)}g(u,y)
  +\frac{4g(x,u)g(y,u)}{g(u,u)^2}g(u,u)\\
  &=g(x,y).
\end{aligned}
\]
因此 $S_u$ 是正交变换. 显然有 $S_u(u)=-u$. 又若 $u\perp v$，则 $S_u(v)=v$. $S_u$ 称为 $V$ 上的镜像变换，它把 $u$ 映为 $-u$，而一切与 $u$ 正交的向量保持不动. 由此不难验证 $S_u^2=I$，即 $S_u$ 是一个对合变换.

由 §9.4 习题 12 知道，$n$ 维欧氏空间中任一正交变换均可表示为不超过 $n$ 个镜像变换之积. 在正交几何中，上述结论依然成立，这就是下面的 Cartan-Dieudonné 定理. 证明这个定理需要进一步引入正交几何学中的一些概念和结论，因此在这里我们就不给出相应的证明了，有兴趣的读者可以参考 [4].

\noindent\textbf{定理 10.5.2}\quad 设 $V$ 是数域 $\mathbb K$ 上的 $n$ 维线性空间，$g$ 是 $V$ 上非退化的对称型. 若 $\eta$ 是 $V$ 上的正交变换，则 $\eta$ 可以表示为不超过 $n$ 个镜像变换之积.

正交几何有许多与欧氏几何相仿的性质，但也有一个明显的不同. 在欧氏空间中，任一非零向量不能与自身垂直，但在正交几何中这是可能的.

\noindent\textbf{定义 10.5.2}\quad 设 $V$ 是有限维线性空间，$g$ 是 $V$ 上非退化的对称型. 若 $v$ 是 $V$ 中的非零向量，且 $g(v,v)=0$，则称 $v$ 是 $V$ 上的一个迷向向量. 含有迷向向量的子空间称为迷向子空间，不含任何迷向向量的子空间称为全不迷向子空间. 若一个子空间的非零向量全是迷向向量，则称之为全迷向子空间.

一个二维线性空间若带有一个非退化的对称型且含有迷向向量，则称之为双曲平面. 对双曲平面有如下基本结论.

\noindent\textbf{定理 10.5.3}\quad (1) $V$ 是双曲平面的充分必要条件是 $V$ 有一组基 $\{u,v\}$，满足条件：
\[
  g(u,u)=g(v,v)=0,
\]
\[
  g(u,v)=g(v,u)=1,
\]
即 $g$ 在这组基下的表示矩阵为
\[
  \left(\begin{array}{cc}
  0&1\\
  1&0
  \end{array}\right);
\]

(2) 任何两个双曲平面皆保距同构；

(3) 任何一个双曲平面有且只有两个一维的全迷向子空间.

\noindent\textbf{证明}\quad (1) 设 $V$ 是双曲平面，则存在 $u\ne0$，使 $g(u,u)=0$. 因为 $g$ 是非退化的，故存在 $v_0\in V$，使 $g(u,v_0)\ne0$. 显然 $u,v_0$ 线性无关，故 $\{u,v_0\}$ 构成 $V$ 的一组基. 不失一般性，可令 $g(u,v_0)=1$（否则在 $v_0$ 上乘一个常数）. 又若 $k\in\mathbb K$，则
\[
  g(ku+v_0,ku+v_0)=g(v_0,v_0)+2k,
\]
令 $v=v_0-\dfrac12g(v_0,v_0)u$，便有 $g(v,v)=0$，而这时仍有 $g(u,v)=1$. 因此我们可以找到所需要的基. 逆命题是显然的.

(2) 由 (1) 可设两个双曲平面分别有一组基 $\{u,v\}$，$\{u',v'\}$ 适合 (1) 中的条件，作线性映射 $\varphi$，使 $\varphi(u)=u'$，$\varphi(v)=v'$. 不难验证 $\varphi$ 是保距同构.

(3) 设 $\{u,v\}$ 是适合 (1) 中条件的一组基. 对任意的 $k,t\in\mathbb K$，
\[
  g(ku+tv,ku+tv)=2kt,
\]
因此 $ku+tv$ 为迷向向量的充分必要条件是 $k\ne0,\ t=0$ 或 $k=0,\ t\ne0$. 这就是说，分别由 $u,v$ 生成的子空间是 $V$ 仅有的两个全迷向子空间. $\square$

\noindent\textbf{例 10.5.1}\quad 设 $V$ 是实数域上的四维空间，若 $g$ 是一个非退化的对称型且其正惯性指数等于 3，则称 $(V,g)$ 是一个 Minkowski（闵可夫斯基）空间. $g$ 在 $V$ 适当的一组基下的表示矩阵为
\[
  \left(\begin{array}{rrrr}
  1&0&0&0\\
  0&1&0&0\\
  0&0&1&0\\
  0&0&0&-1
  \end{array}\right).
\]
$V$ 上的正交变换称为 Lorentz（洛伦兹）变换，$V$ 中的迷向向量称为光向量，$V$ 中适合 $g(x,x)>0$ 的向量 $x$ 称为空间向量，而适合 $g(x,x)<0$ 的向量 $x$ 称为时间向量. Minkowski 空间在相对论中有重要的应用.

\section*{习题 10.5}

1. 设 $g$ 是 $V$ 上非退化的对称型，$\varphi$ 是 $V$ 上的正交变换，求证：$\det\varphi=\pm1$.

2. 若 $\varphi$ 是 $(V,g)$ 上的正交变换且 $\det\varphi=1$，求证：$\varphi$ 可以表示为偶数个镜像变换之积.

3. 设 $\varphi$ 是 $(V,g)$ 上的线性变换，$\varphi^*$ 是 $\varphi$ 关于 $g$ 的对偶映射，证明：$\varphi$ 是正交变换的充分必要条件是 $\varphi^*\varphi=I$.

4. 设 $V$ 是双曲平面，$\varphi$ 是 $V$ 上的正交变换且 $\det\varphi=-1$，求证：$\varphi$ 必是镜像变换.

5. 若 $V=U_1\oplus\cdots\oplus U_r$ 且 $U_i\perp U_j(i\ne j)$，则称 $V$ 是 $U_i$ 的正交直和，记为
\[
  V=U_1\perp\cdots\perp U_r.
\]
若 $V=U_1\perp\cdots\perp U_r=V_1\perp\cdots\perp V_r$ 且存在保距同构 $\varphi_i:U_i\to V_i(i=1,\cdots,r)$，求证：存在 $V$ 上的正交变换 $\varphi$，使 $\varphi$ 限制在每个 $U_i$ 上就是 $\varphi_i$.

6. 若 $\varphi$ 是 $(V,g)$ 的正交变换且 $V_1=\{x\in V\mid \varphi(x)=x\}$，求证：
\[
  \dim V=\dim V_1+\dim(I-\varphi)V,
\]

\section*{历史与展望}

1868 年，Weierstrass 完成了二次型的理论并将其推广到双线性型
\[
  a_{11}x_1y_1+a_{12}x_1y_2+\cdots+a_{nn}x_ny_n.
\]
Frobenius 在 1878 年的论文《关于线性替换和双线性型》中，用双线性型的语言极大地发展了矩阵理论，他说：“双线性型可以看成是排成 $n$ 行 $n$ 列的 $n^2$ 个变量的系统.” Frobenius 受其老师 Weierstrass 的启发，并在论文中遵循 Weierstrass 的传统，那就是强调严谨的方法，寻求理论中的基本思想. 例如，他深入研究了双线性型的标准型问题，并将一些特例归功于 Kronecker 和 Weierstrass.

$n$ 维欧氏空间就是一个 $n$ 维实线性空间 $V$ 并带有一个正定对称的双线性型（即内积）$g(x,y)=\sum_{i,j=1}^n a_{ij}x_iy_j$. 如果考虑一般的非退化对称型 $g(x,y)$，就得到了正交几何，它在相对论中有重要的应用. 如果考虑一般的非退化交错型 $g(x,y)$，就得到了辛几何，它在复几何和弦理论中有重要的应用.

\section*{复习题十}

1. 设 $V$ 是数域 $\mathbb K$ 上的线性空间（不必假设维数有限），$f,g$ 是 $V$ 上的非零线性函数，求证：$f$ 和 $g$ 线性相关的充分必要条件是 $\operatorname{Ker} f=\operatorname{Ker} g$.

2. 设 $V$ 是数域 $\mathbb K$ 上的 $n$ 维线性空间，$U$ 是 $V$ 的非平凡子空间，求证：必存在 $V$ 上的线性函数 $f_i(i=1,\cdots,r)$，使
\[
  U=\bigcap_{i=1}^r \operatorname{Ker} f_i.
\]

3. 设 $U,V$ 是数域 $\mathbb K$ 上的有限维线性空间，$U^*,V^*$ 分别是它们的对偶空间. 求证：
\[
  U^*\oplus V^*\cong (U\oplus V)^*.
\]

4. 设 $\varphi$ 是线性空间 $V$（不要求是有限维）上的幂等线性变换（即 $\varphi^2=\varphi$），$\varphi^*$ 是 $\varphi$ 的对偶变换，求证：
\[
  \operatorname{Im}\varphi^*=(\operatorname{Ker}\varphi)^\perp.
\]

5. 设 $V$ 是 $n$ 维欧氏空间，则对任一固定的 $u\in V$，$(u,-)$ 是 $V$ 上的线性函数，作映射 $\eta:V\to V^*,\ \eta(u)=(u,-)$. 证明：

(1) $\eta$ 是线性同构，特别，若将 $u$ 与 $(u,-)$ 等同起来，则 $\langle u,v\rangle=(u,v)$，即可将 $V$ 看成是自身的对偶空间；

(2) $V$ 的任一组标准正交基 $e_1,e_2,\cdots,e_n$ 的对偶基是其自身；

(3) $V$ 上任一线性变换 $\varphi$ 的对偶变换就是 $\varphi$ 的伴随.

6. 设 $U,V$ 是 $\mathbb K$ 上的线性空间，$g$ 是 $U$ 和 $V$ 上的非退化双线性型，$\varphi,\psi$ 是 $V$ 上的线性变换. 证明：

(1) $(k\varphi+l\psi)^*=k\varphi^*+l\psi^*$，其中 $k,l\in\mathbb K$；

(2) $(\psi\varphi)^*=\varphi^*\psi^*$；

(3) 若 $\varphi$ 是 $V$ 的自同构，则 $\varphi^*$ 是 $U$ 的自同构，此时 $(\varphi^*)^{-1}=(\varphi^{-1})^*$；

(4) $(\varphi^*)^*=\varphi$.

7. 设 $h$ 是三维线性空间 $V$ 上的非零交错型，求证：存在 $V$ 上的线性函数 $f,g$，使对任意的 $x,y\in V$，有
\[
  h(x,y)=f(x)g(y)-f(y)g(x).
\]

8. 设 $g$ 是 $n$ 维实线性空间 $V$ 上的非退化对称型，$g$ 的正惯性指数为 $p$，负惯性指数为 $q$. 设 $W$ 是 $V$ 的极大全迷向子空间，求证：
\[
  \dim W=\min\{p,q\}.
\]

9. 证明 Minkowski 空间的下列性质：

(1) 任意两个时间向量不可能互相正交；

(2) 任意一个时间向量不可能正交于一个光向量；

(3) 两个光向量正交的充分必要条件是它们线性相关.

\end{document}
